\setcounter{section}{5}

  \zwischenueberschrift{Soal latihan}

\inputaufgabegibtloesung
{}
{

Misalkan

\mavergleichskettedisp
{\vergleichskette
{ f(x,y)
}
{ =} { ax^2+by^2+cxy
}
{ } {
}
{ } {
}
{ } {
}
}
{}{}{.}
Tentukan
\definitionsverweis {kelengkungan utama}{}{}
dan
\definitionsverweis {arah kelengkungan utama}{}{}
pada
\definitionsverweis {grafik}{}{}
fungsi $f$ di titik nol.

}
{} {}

\inputaufgabe
{}
{

Tentukan
\definitionsverweis {kelengkungan utama}{}{}
dan
\definitionsverweis {arah kelengkungan utama}{}{}
pada setiap titik
\definitionsverweis {grafik}{}{}
fungsi
\maabbeledisp {f} {\R^{n-1}} {\R
} { \left(  x_1   , \,   \ldots  , \,   x_{n-1}                 \right) } { a_1x_1^2  + \cdots + a_{n-1} x_{n-1}^2
} {.}

}
{} {}

\inputaufgabe
{}
{

Misalkan

\mavergleichskettedisp
{\vergleichskette
{ f(x,y,z)
}
{ =} { ax^2 +by^2 +c z^2 +dxy
}
{ } {
}
{ } {
}
{ } {
}
}
{}{}{.}
Hitung
\definitionsverweis {polinom karakteristik}{}{}
dari
\definitionsverweis {pemetaan Weingarten}{}{}
pada
\definitionsverweis {grafik}{}{}
fungsi $f$ di titik nol. Hitung pula
\definitionsverweis {kelengkungan utama}{}{}
dan
\definitionsverweis {arah kelengkungan utama}{}{.}

}
{} {}

\inputaufgabe
{}
{

Misalkan

\mavergleichskettedisp
{\vergleichskette
{ f(x,y,z)
}
{ =} { ax^2 +by^2 +c z^2 +dxy + e yz
}
{ } {
}
{ } {
}
{ } {
}
}
{}{}{.}
Hitung
\definitionsverweis {polinom karakteristik}{}{}
dari
\definitionsverweis {pemetaan Weingarten}{}{}
pada
\definitionsverweis {grafik}{}{}
fungsi $f$ di titik nol. Dalam kasus ini, seberapa sulitkah menentukan
\definitionsverweis {kelengkungan utama}{}{?}

}
{} {}

\inputaufgabe
{}
{

Misalkan

\mavergleichskettedisp
{\vergleichskette
{ f(x,y)
}
{ =} { x^2+y^3
}
{ } {
}
{ } {
}
{ } {
}
}
{}{}{.}
Tentukan
\definitionsverweis {kelengkungan utama}{}{,}
\definitionsverweis {arah kelengkungan utama}{}{}
dan
\definitionsverweis {kelengkungan Gauss}{}{}
pada
\definitionsverweis {grafik}{}{}
fungsi $f$ di setiap titik

\mavergleichskette
{\vergleichskette
{ P
}
{ = }{ (Q,f(Q))
}
{  }{
}
{    }{
}
{   }{
}
}
{}{}{,}

\mavergleichskette
{\vergleichskette
{ Q
}
{ \in }{ \R^2
}
{  }{
}
{    }{
}
{   }{
}
}
{}{}{.}

}
{} {}

\inputaufgabe
{}
{

Misalkan

\mavergleichskette
{\vergleichskette
{ V
}
{ \subseteq }{ \R^{n-1}
}
{  }{
}
{    }{
}
{   }{
}
}
{}{}{}
\definitionsverweis {terbuka}{}{}
dan misalkan
\maabb {f} { V } {\R
} {}
adalah suatu
\definitionsverweis {fungsi yang dua kali terdiferensialkan secara kontinu}{}{.}
Buktikan bahwa
\definitionsverweis {kelengkungan utama}{}{}
dan
\definitionsverweis {arah kelengkungan utama}{}{}
pada
\definitionsverweis {grafik}{}{}
fungsi $f$ di suatu titik

\mavergleichskette
{\vergleichskette
{ P
}
{ = }{ (Q,f(Q))
}
{  }{
}
{    }{
}
{   }{
}
}
{}{}{}
hanya bergantung pada turunan-turunan parsial pertama dan kedua dari $f$ di $Q$.

}
{} {}

\inputaufgabe
{}
{

Misalkan

\mavergleichskette
{\vergleichskette
{ W
}
{ \subseteq }{ \R^n
}
{  }{
}
{    }{
}
{   }{
}
}
{}{}{}
\definitionsverweis {terbuka}{}{,}
\maabb {h} { W } { \R
} {}
adalah suatu
\definitionsverweis {fungsi yang dua kali terdiferensialkan secara kontinu}{}{}
dan

\mavergleichskette
{\vergleichskette
{ Y
}
{ = }{ h^{-1}(c)
}
{  }{
}
{    }{
}
{   }{
}
}
{}{}{}
adalah
\definitionsverweis {serat}{}{}
di atas

\mavergleichskette
{\vergleichskette
{ c
}
{ \in }{ \R
}
{  }{
}
{    }{
}
{   }{
}
}
{}{}{,}
dengan $h$
\definitionsverweis {reguler}{}{}
di setiap titik pada $Y$.
Catatan edisi: sumber hanya mengasumsikan $h$ kelas $C^1$, padahal kelengkungan utama dan arah utamanya memerlukan pemetaan Weingarten. Hipotesis di atas diperkuat menjadi $C^2$ agar semua objek yang ditanyakan terdefinisi sebagaimana pada kuliah.
Kita juga memandang $h$ sebagai suatu pemetaan $\tilde{h}$ pada
\mathl{W \times \R^m}{},
dengan $m$ variabel terakhir tidak muncul secara eksplisit dalam $\tilde{h}$. Misalkan

\mavergleichskettedisp
{\vergleichskette
{ Z
}
{ =} { \tilde{h}^{-1}(c)
}
{ =} { Y \times \R^m
}
{ \subseteq} { W \times \R^{m}
}
{ } {
}
}
{}{}{.}
Misalkan

\mavergleichskette
{\vergleichskette
{ Q
}
{ \in }{ Z
}
{  }{
}
{    }{
}
{   }{
}
}
{}{}{}
adalah suatu titik yang terletak di atas

\mavergleichskette
{\vergleichskette
{ P
}
{ \in }{ Y
}
{  }{
}
{    }{
}
{   }{
}
}
{}{}{.}
Bagaimana hubungan antara
\definitionsverweis {kelengkungan utama}{}{}
dan
\definitionsverweis {arah kelengkungan utama}{}{}
dari $Z$ di $Q$ dengan yang dari $Y$ di $P$?

}
{} {}

Misalkan

\mavergleichskette
{\vergleichskette
{ W
}
{ \subseteq }{ \R^n
}
{  }{
}
{    }{
}
{   }{
}
}
{}{}{}
\definitionsverweis {terbuka}{}{,}
\maabb {h} { W } { \R
} {}
adalah suatu
\definitionsverweis {fungsi yang dua kali terdiferensialkan secara kontinu}{}{}
dan

\mavergleichskette
{\vergleichskette
{ Y
}
{ = }{ h^{-1}(c)
}
{  }{
}
{    }{
}
{   }{
}
}
{}{}{}
adalah
\definitionsverweis {serat}{}{}
di atas

\mavergleichskette
{\vergleichskette
{ c
}
{ \in }{ \R
}
{  }{
}
{    }{
}
{   }{
}
}
{}{}{,}
dengan $h$
\definitionsverweis {reguler}{}{}
di setiap titik pada $Y$. Untuk suatu titik

\mavergleichskette
{\vergleichskette
{ P
}
{ \in }{ Y
}
{  }{
}
{    }{
}
{   }{
}
}
{}{}{}
\definitionsverweis {determinan}{}{}
dari
\definitionsverweis {pemetaan Weingarten}{}{}
$L_P$ disebut
\definitionswort {kelengkungan Gauss--Kronecker}{}
dari $Y$ di $P$.

Kelengkungan Gauss--Kronecker merupakan perumuman langsung dari kelengkungan Gauss.

\inputaufgabe
{}
{

Misalkan

\mavergleichskette
{\vergleichskette
{ W
}
{ \subseteq }{ \R^n
}
{  }{
}
{    }{
}
{   }{
}
}
{}{}{}
\definitionsverweis {terbuka}{}{,}
\maabb {h} { W } { \R
} {}
adalah suatu
\definitionsverweis {fungsi yang dua kali terdiferensialkan secara kontinu}{}{}
dan

\mavergleichskette
{\vergleichskette
{ Y
}
{ = }{ h^{-1}(c)
}
{  }{
}
{    }{
}
{   }{
}
}
{}{}{}
adalah
\definitionsverweis {serat}{}{}
di atas

\mavergleichskette
{\vergleichskette
{ c
}
{ \in }{ \R
}
{  }{
}
{    }{
}
{   }{
}
}
{}{}{,}
dengan $h$
\definitionsverweis {reguler}{}{}
di setiap titik pada $Y$. Misalkan

\mavergleichskette
{\vergleichskette
{ P
}
{ \in }{ Y
}
{  }{
}
{    }{
}
{   }{
}
}
{}{}{.}
Buktikan bahwa hasil kali
\definitionsverweis {kelengkungan utama}{}{}
dari $Y$ di $P$ sama dengan
\definitionsverweis {kelengkungan Gauss--Kronecker}{}{}
dari $Y$ di $P$.

}
{} {}

\inputaufgabe
{}
{

Misalkan

\mavergleichskettedisp
{\vergleichskette
{ h(x,y,z,w)
}
{ =} { x^2+y^3+z^4+w^5
}
{ } {
}
{ } {
}
{ } {
}
}
{}{}{}
dan

\mavergleichskette
{\vergleichskette
{ Y
}
{ = }{ h^{-1}(1)
}
{  }{
}
{    }{
}
{   }{
}
}
{}{}{.}
Untuk titik

\mavergleichskettedisp
{\vergleichskette
{ P
}
{ =} { \left(  1   , \,   -1  , \,   1 , \,   0                \right)
}
{ \in} { Y
}
{ } {
}
{ } {
}
}
{}{}{}
tentukan
\definitionsverweis {pemetaan Weingarten}{}{}
$L_P$,
\definitionsverweis {polinom karakteristik}{}{}
darinya, serta
\definitionsverweis {kelengkungan Gauss--Kronecker}{}{}
dari $Y$ di $P$.

}
{} {}

\inputaufgabe
{}
{

Misalkan

\mavergleichskettedisp
{\vergleichskette
{ h(x,y,z)
}
{ =} { x^2+y^3+z^5
}
{ } {
}
{ } {
}
{ } {
}
}
{}{}{}
pada

\mavergleichskette
{\vergleichskette
{ W
}
{ = }{ \R^3 \setminus \{0\}
}
{  }{
}
{    }{
}
{   }{
}
}
{}{}{}
dan

\mavergleichskette
{\vergleichskette
{ Y
}
{ = }{ h^{-1}(0)
}
{  }{
}
{    }{
}
{   }{
}
}
{}{}{.}
Tentukan
\definitionsverweis {kelengkungan normal}{}{}
di titik

\mavergleichskette
{\vergleichskette
{ P
}
{ = }{ \begin{pmatrix}  1  \\ -1\\  0   \end{pmatrix}
}
{  }{
}
{    }{
}
{   }{
}
}
{}{}{}
dalam arah vektor singgung ternormalisasi
\mathl{{ \frac{   1  }{  \sqrt{29}  } }  \begin{pmatrix}  3  \\ -2\\  4   \end{pmatrix}}{.}

}
{} {}

  \zwischenueberschrift{Soal untuk dikumpulkan}

\inputaufgabe
{4}
{

Misalkan

\mavergleichskettedisp
{\vergleichskette
{ f(x,y)
}
{ =} { x^3 +x^2y+3xy^2-y^4
}
{ } {
}
{ } {
}
{ } {
}
}
{}{}{.}
Tentukan
\definitionsverweis {kelengkungan utama}{}{,}
\definitionsverweis {arah kelengkungan utama}{}{}
dan
\definitionsverweis {kelengkungan Gauss}{}{}
pada
\definitionsverweis {grafik}{}{}
fungsi $f$ di setiap titik

\mavergleichskette
{\vergleichskette
{ P
}
{ = }{ (Q,f(Q))
}
{  }{
}
{    }{
}
{   }{
}
}
{}{}{,}

\mavergleichskette
{\vergleichskette
{ Q
}
{ \in }{ \R^2
}
{  }{
}
{    }{
}
{   }{
}
}
{}{}{.}

}
{} {}

\inputaufgabe
{4}
{

Misalkan

\mavergleichskettedisp
{\vergleichskette
{ f(x,y,z)
}
{ =} { -7x^2 +3y^2 -5z^2+ 6xy
}
{ } {
}
{ } {
}
{ } {
}
}
{}{}{.}
Tentukan
\definitionsverweis {kelengkungan utama}{}{}
dan
\definitionsverweis {arah kelengkungan utama}{}{}
pada
\definitionsverweis {grafik}{}{}
fungsi $f$ di titik nol.

}
{} {}

\inputaufgabe
{6}
{

Misalkan

\mavergleichskettedisp
{\vergleichskette
{ h(x,y,z,w)
}
{ =} { x^2+yz+z^3+xyzw
}
{ } {
}
{ } {
}
{ } {
}
}
{}{}{}
dan

\mavergleichskette
{\vergleichskette
{ Y
}
{ = }{ h^{-1}(2)
}
{  }{
}
{    }{
}
{   }{
}
}
{}{}{.}
Untuk titik

\mavergleichskettedisp
{\vergleichskette
{ P
}
{ =} { \left(  1   , \,   1  , \,   1 , \,   -1                \right)
}
{ \in} { Y
}
{ } {
}
{ } {
}
}
{}{}{}
tentukan
\definitionsverweis {pemetaan Weingarten}{}{}
$L_P$,
\definitionsverweis {polinom karakteristik}{}{}
darinya, serta
\definitionsverweis {kelengkungan Gauss--Kronecker}{}{}
dari $Y$ di $P$.

}
{} {Petunjuk: Selesaikan persamaan terhadap suatu variabel yang sesuai dan perlakukan hipermuka tersebut sebagai grafik dalam variabel-variabel lainnya.}

\inputaufgabe
{6 (2+2+2)}
{

Misalkan

\mavergleichskettedisp
{\vergleichskette
{ h(x,y,z)
}
{ =} { x^3+y^3+z^3
}
{ } {
}
{ } {
}
{ } {
}
}
{}{}{}
pada

\mavergleichskette
{\vergleichskette
{ W
}
{ = }{ \R^3 \setminus \{0\}
}
{  }{
}
{    }{
}
{   }{
}
}
{}{}{}
dan

\mavergleichskette
{\vergleichskette
{ Y
}
{ = }{ h^{-1}(0)
}
{  }{
}
{    }{
}
{   }{
}
}
{}{}{.}
Misalkan

\mavergleichskette
{\vergleichskette
{ P
}
{ = }{ \begin{pmatrix}  1  \\ 1\\  - \sqrt[3]{2}   \end{pmatrix}
}
{ \in }{ Y
}
{    }{
}
{   }{
}
}
{}{}{}
dan

\mavergleichskettedisp
{\vergleichskette
{ v
}
{ =} { \begin{pmatrix}  1  \\ -1\\  0   \end{pmatrix}
}
{ \in} { T_PY
}
{ } {
}
{ } {
}
}
{}{}{.}
\aufzaehlungdrei{Tentukan suatu persamaan bidang untuk
\definitionsverweis {bidang normal}{}{}
yang bersesuaian dengan $v$.
}{Tentukan
\definitionsverweis {kelengkungan normal}{}{}
di $P$ dalam arah vektor singgung ternormalisasi
\mathl{{ \frac{   v  }{  \Vert { v } \Vert  } }}{.}
}{Tentukan suatu kurva berparametrisasi panjang busur
\maabbdisp {\gamma} { I } { Y
} {}
dengan

\mavergleichskette
{\vergleichskette
{ \gamma (0)
}
{ = }{ P
}
{  }{
}
{    }{
}
{   }{
}
}
{}{}{,}

\mavergleichskette
{\vergleichskette
{ \gamma' (0)
}
{ = }{ { \frac{   v  }{  \Vert { v } \Vert  } }
}
{  }{
}
{    }{
}
{   }{
}
}
{}{}{.}
Verifikasikan berlakunya
Teorema 5.14
dalam situasi ini.
}

Catatan edisi: pada butir ketiga, orientasikan bidang normal dengan menetapkan pasangan $(v/\lVert v\rVert,N(P))$ sebagai basis positif, sehingga normal satuan kurva irisan di dalam bidang yang bersesuaian dengan arah singgung $v/\lVert v\rVert$ adalah $N(P)$. Tanpa konvensi ini, verifikasi kelengkungan bertanda hanya menentukan kesamaan nilai mutlak.

}
{} {}

\inputaufgabe
{2}
{

Buatlah sketsa suatu permukaan terdiferensialkan $Y$ di ruang dan suatu bidang normal $E$ yang irisannya dengan permukaan tersebut menghasilkan suatu kurva yang tidak reguler pada setidaknya satu titik.

}
{} {}
